\section{Related Work}
\label{sec:related}

\wcabench{} sits at the intersection of five research lines: \wca{} data
analysis, cube-solving benchmarks, sports analytics, time-series benchmarking,
and the methodology of dataset/benchmark papers.

\subsection{Speedcubing and World Cube Association data}
\label{sec:related-wca}

The \wca{} maintains the world's largest structured record of speedcubing
results, published as a versioned database export and governed by an explicit
rule book \citep{wca_export, wca_regs}. Prior work on this data has
concentrated on three narrow questions. Placement has been analysed with
dynamic paired-comparison and density-based rating models
\citep{glickman_rating}. World-record progression has been studied with
regression and saturating-curve models \citep{tatem_sprint, nevill_running}.
Human limits have been approached with extreme-value theory and saturating
curves \citep{evt_coles, nevill_running}. Each of these studies optimizes a
single task under its own protocol. Our contribution is not a better estimator
for any one of them but a common protocol under which all of them---and new
methods---can be evaluated and compared.

\subsection{Cube-solving benchmarks}
\label{sec:related-cubebench}

A parallel literature evaluates whether machines can \emph{solve} cubes. Early
work learned to solve the three-by-three cube with deep reinforcement learning
and search \citep{deepcubea} or with approximate policy iteration
\citep{mcaleer_cube}. More recently, a family of \emph{cube-reasoning
benchmarks} has emerged. CubeBench \citep{cubebench} builds a three-tiered
generative benchmark on the Rubik's cube to diagnose three cognitive challenges
that hinder LLM agents in physically grounded settings---spatial reasoning,
long-horizon state tracking, and active exploration under partial observations.
Cube Bench \citep{cube_bench_mllm} evaluates multimodal LLMs on five
spatial/sequential skills, from reconstructing cube faces from images to
multi-step planning and self-correction.

These benchmarks and ours are \emph{complementary rather than competing}: the
cube-reasoning benchmarks measure an agent's reasoning and planning ability on
\emph{synthetic} cube states, whereas \wcabench{} measures the predictive and
inferential ability of statistical and machine-learning models on \emph{real,
noisy, rule-governed competition records}. We therefore adopt a different
evaluation object (statistical/ML models rather than LLM/MLLM agents), a
different data source (real competitions rather than generated scrambles) and a
different task family (regression, ranking, classification, extreme-value
estimation and causal inference).

\subsection{Sports analytics}
\label{sec:related-sports}

Sports analytics has a long tradition of forecasting and rating. In
individual sports, performance-progression models for running and swimming
records are closely related to our \task{4} \citep{tatem_sprint,
nevill_running}. In team sports, ranking and outcome models built on
Bradley--Terry and Elo-style priors \citep{bradley_terry, elo} are the
ancestors of our \task{2} baselines, with the Plackett--Luce model
\citep{plackett_luce} providing the probabilistic ranking generalization we
use. In cricket, the Duckworth--Lewis method and its successors quantify how
context changes performance targets, an example of a \emph{rule-aware}
performance model. Injury- and fatigue-risk modelling in soccer and basketball
\citep{injury_prediction} resembles our \task{3} in that it predicts rare,
consequential events from longitudinal records. Our benchmark borrows the
discipline of these fields---time-aware splits, calibration, effect sizes---and
packages it for a new domain.

\subsection{Time-series and tabular benchmarking}
\label{sec:related-ts}

Standardized forecasting competitions and corpora have shaped evaluation
practice, notably the M-competitions and the M5 accuracy competition
\citep{m5_competition}, the Monash Time Series Forecasting Archive
\citep{monash_tsf} and the GluonTS toolkit \citep{gluonts}. These resources
established the norms we follow: a fixed, frozen test window; explicit
avoidance of future leakage; stratified reporting; and public leaderboards,
building on standard forecasting methodology \citep{hyndman_fpp3}.
\wcabench{} differs in that its series are indexed by irregular competition
calendars rather than regular timestamps, and in that the prediction targets
are governed by formal round rules. For online/incremental learning, the
prequential (test-then-train) evaluation tradition \citep{prequential} is the
closest methodological relative of our rolling-window protocol.

\subsection{Extreme-value estimation and human limits}
\label{sec:related-evt}

Estimating an unobservable performance ceiling is an extrapolation problem.
Extreme-value theory provides the asymptotic machinery for tail behaviour
\citep{evt_coles, pickands}, while saturating-curve models are the pragmatic
choice in sports \citep{nevill_running, tatem_sprint}. Bayesian hierarchical models let
sports with scarce data borrow strength from a population of similar events
\citep{bayesian_hierarchical}. Our \task{4} benchmark exposes exactly this
tension: recurring events have dense histories, while emerging or discontinued
events may have fewer than fifty world-record points, forcing explicit
modelling choices that we make comparable across methods.

\subsection{Causal inference in observational sports data}
\label{sec:related-causal}

Our \task{5} formalizes a question---does practising event $A$ causally improve
event $B$?---that is deceptively hard because participation is self-selected.
Difference-in-differences \citep{did}, instrumental variables \citep{iv},
regression discontinuity, and causal forests \citep{causal_forest} are the
standard tools, each resting on a different identification assumption. This
mirrors causal inference in economics and epidemiology, where the honest
reporting of assumptions and sensitivity analyses has become standard
\citep{sensitivity_analysis}. We deliberately include a naive correlational
baseline to make the gap between association and causation measurable.

\subsection{Benchmark and dataset methodology}
\label{sec:related-methodology}

Our documentation and reporting follow the datasheet \citep{datasheets} and
model-card \citep{model_cards} traditions, which argue that datasets and
models must document provenance, composition, recommended uses and limitations.
We adopt the NeurIPS Datasets and Benchmarks evaluation philosophy
\citep{neurips_dnb} by requiring stratified results, calibration and effect
sizes rather than a single headline number. Our leakage-aware protocol echoes
broader warnings about spurious leakage in machine learning evaluation
\citep{leakage_kaufman}.

\paragraph{Deep and graph methods in this setting.} Our measured multi-seed
results speak directly to the above lines of work
(Section~\ref{sec:results-multiseed}): on \task{1} a sequence model
(\maelog{} $0.9436\pm0.0754$) is roughly nine times worse than tree boosting
($0.1055\pm0.0417$), and on \task{2} a graph neural network
($\kendall=0.7440\pm0.1258$) trails a simple domain ranking
($0.8132\pm0.0342$). In rule-governed longitudinal sports records, domain-aware
features currently appear to matter more than generic deep or graph
architectures.
